% ═══════════════════════════════════════════════════════════════════════
\section{Synthèse}
% ═══════════════════════════════════════════════════════════════════════

\begin{frame}{Ce qu'il faut retenir}
  \small
  \begin{itemize}
    \item Une métrique : quatre axiomes (positivité,
      séparation, symétrie, inégalité triangulaire) --- sur
      n'importe quel ensemble.
    \item Un même ensemble peut porter plusieurs métriques ($\dE$,
      $\dT$) ; la forme des boules en dépend.
    \item Des métriques « appliquées » : distance sphérique
      (navigation), distance de Hamming (codes correcteurs).
    \item Ouvert, fermé, compact ; dans $\R^n$,
      compact = fermé et borné (Heine-Borel).
    \item La distance de Hausdorff fait de $\Comp{E}$ un espace
      métrique dont les points sont des ensembles.
    \item Complétude : toute suite de Cauchy converge dans
      l'espace ($\R^n$ oui, $\Q$ non).
    \item Banach : dans un espace complet, une contraction a un
      unique point fixe, atteint en itérant depuis n'importe où, avec
      une erreur en $k^n/(1-k)$.
  \end{itemize}
\end{frame}

\begin{frame}{Le théorème de Banach, une dernière fois}
  \FigBox[0.36\textheight]{%
    \begin{tikzpicture}[
        hyp/.style={draw=sectioncolor, fill=thmboxbg, rounded corners=2pt,
          thick, align=center, font=\small, minimum height=1.1cm,
          text width=3.2cm},
        conc/.style={hyp, draw=coveraccent!85!black, fill=secboxbg,
          text width=4.4cm}]
      \node[hyp] (e) at (0,0.75) {$(E,d)$ complet};
      \node[hyp] (f) at (0,-0.75) {$f$ contraction\\$k<1$};
      \node[conc] (c) at (6.2,0) {un unique $\overline{x}$\\
        $f^n(x_0)\to\overline{x}$ pour tout $x_0$\\
        erreur $\leq\frac{k^n}{1-k}\,d(x_0,x_1)$};
      \draw[fleche] (e.east) -- (c.west);
      \draw[fleche] (f.east) -- (c.west);
    \end{tikzpicture}}
  \small
  \begin{center}
  \begin{tabular}{ll}
    \toprule
    Espace complet & Ce qu'on y trouve comme point fixe \\
    \midrule
    $I\subset\R$ fermé & le zéro cherché par la méthode de Newton \\
    $\big(\Comp{\R^2},\dH\big)$ & une fractale (chapitre 3) \\
    vecteurs de probabilité & le classement PageRank (chapitre 4) \\
    \bottomrule
  \end{tabular}
  \end{center}
\end{frame}
